% !TEX root = scientific-computing.tex

\chapter{Parallel Computing} \label{ch:parallel}

Briefly, parallel computing is a method of running code on multiple
processors (or multiple cores of the same processor) at the same time.
In general, this is a difficult task depending on where data is stored
and retrieved. The
\href{https://docs.julialang.org/en/v1/manual/parallel-computing/}{Julia
Documentation on parallel computing} is a good place to start.

Let's start with something relatively simple. Consider the following
code:

\begin{jlcode}[parallel]
 function countHeads(n::Int)
    c::Int = 0
    for i=1:n
        c += rand(Bool)
    end
    c
end
\end{jlcode}

which mimics flipping a coin \texttt{n} times. We can simulate flipping
2 billion coins and finding the fraction that is heads by the
following:

\begin{jlverbatim}
@time countHeads(2*10^9)/(2*10^9)
\end{jlverbatim}
%
which returns \jlc[parallel]{display(@time countHeads(2*10^9)/(2*10^9))}\printpythontex[verbatim]


\section{Running parallel code}

We now wish to take advantage of the fact that today's processors often
have multiple cores and multiple threading. There are some helpful
things in the \texttt{Distributed} package:

\begin{jlblock}[parallel]
using Distributed
\end{jlblock}

Let's say that we wish to run it simultaneously on 2 cores (even if we
know we have more) If we type

\begin{jlverbatim}
addproc(1)
\end{jlverbatim}

this would then allow julia up to 2 ``processors''. There are a couple of ways to now run code. The first gives us more control, while the second allows for some easy code writing.


\subsection{Assigning code to workers}

First, we need to make sure that the function \verb!countHeads! is available to each core of the machine. To do this, we will start the function with the \texttt{@everywhere} macro:

\begin{jlblock}[parallel]
@everywhere function countHeads(n::Int)
   c::Int = 0
   for i=1:n
       c += rand(Bool)
   end
   c
end
\end{jlblock}

Any function that you will need on multiple cores will need to be
prefaced with the macro \texttt{@everywhere}. After this, we can now send code to different cores. Type:
%
\begin{jlblock}[parallel]
a= @spawn countHeads(10^9);
b= @spawn countHeads(10^9)
\end{jlblock}
%
and you should notice that it doesn't return the number of heads as
expected and it goes almost instanteously. Don't worry right now about
the type of object that is returned.

Even though it returns quickly, code is already running and you can see
this with your machine's process manager (Activity Monitor on MacOS). To get the results, we type
%
\begin{jlblock}[parallel]
fetch(a)+fetch(b)
\end{jlblock}
%
and it will return the number of heads flipped. To test the timing, put all three of these in a block:

\begin{jlverbatim}[parallel]
a= @spawn countHeads(10^9);
b= @spawn countHeads(10^9);
@time fetch(a)+fetch(b)
\end{jlverbatim}
\begin{jlcode}[parallel]
a= @spawn countHeads(10^9);
b= @spawn countHeads(10^9);
display(@time fetch(a)+fetch(b))
\end{jlcode}
%
and the results are: \printpythontex[verbatim]

This is perhaps a bit disappointing in that if we use two cores, then might expect half of the time.  It's not quite that simple in that there is some overhead for using multiple cores.  However, if you use more cores like the exercise below, you'll probably get better results and for longer processes, it will pay off.  

\subsection{Exercise}

\begin{enumerate}

\item  Add two more cores to julia with the \texttt{addprocs} command.
\item  Rerun the block of code  \verb!@everywhere function countHeads! code above.
\item  Create 4 spawning lines. Call them \texttt{a,\ b,\ c} and \texttt{d}
  and use \verb!5*10^8! for each.
\item  Time the sum of the four.
\item  Note: if you truly have 4 cores, then you will see a further halving of the time. If not, you will probably only see the same results.
\end{enumerate}


\subsection{Adding all of the cores}

Since most computers have multiple cores, this is helpful, but what
about adding all of the cores? If we just call

\begin{jlblock}[parallel]
addprocs()
\end{jlblock}

then it adds everything available. You will get an array of the workers.


\subsection{Parallel for loops}

One issue with what we did above is that we have to think about spawning
to individual processors or cores. That is fairly annoying. Fortunately,
another helpful feature of julia is that of a parallel for loop. Try the
following code:

\begin{jlblock}[parallel]
@time let
 nheads = @distributed (+) for i = 1:2*10^10
   Int(rand(Bool))
 end
end
\end{jlblock}
and the results are:
%\begin{jlcode}[parallel]
%@time display(let
% nheads = @distributed (+) for i = 1:2*10^10
%   Int(rand(Bool))
% end
%end)
%\end{jlcode} 
\printpythontex[verbatim]

where julia distributes the load in the for loop across the available
processors and the time should be a fraction of the original head counter code
where the fraction is the number of actual number of possible processors
or cores.

\section{Writing a parallel card simulator}

For this, we need to reload the PlayingCard module:
%
\begin{jlverbatim}[parallel]
include("PlayingCards.jl")
using .PlayingCards, Random
\end{jlverbatim}
\begin{jlcode}[parallel]
include("../julia-files/PlayingCards.jl")
using .PlayingCards, Random
\end{jlcode}

(or you may need to put in a different path to your module)

If we have the following code to simulate a large number of hands:

\begin{jlblock}[parallel]
function countHands(trials::Int,f::Function)
    local deck=map(Card,1:52)
    local num_hands=0
    for i=1:trials
        shuffle!(deck)
        h = Hand(deck[1:5])
        if(f(h))
            num_hands+=1
        end
    end
    num_hands
end
\end{jlblock}

And let's just check to make sure it works:

\begin{jlblock}[parallel]
@time countHands(10_000_000,isFullHouse)
\end{jlblock} 
which returns \jlc[parallel]{display(@time countHands(10_000_000,isFullHouse))} 
\printpythontex[verbatim]


We now wish to write a distributed version of this in which we will
replace the for loop with a distributed for loop.  First, the module needs to be loaded for all cores by
\begin{jlverbatim}
@everywhere include("PlayingCards.jl")
@everywhere using .PlayingCards, Random
\end{jlverbatim}
\begin{jlcode}[parallel]
@everywhere include("julia-files/PlayingCards.jl")
@everywhere using .PlayingCards, Random
\end{jlcode}
%
and then the following will be the parallel version of the hand count function:
%
%
\begin{jlblock}[parallel]
@everywhere function paraCountHands(trials::Integer,f::Function)
  local deck=map(Card,1:52)
  function checkHand(f::Function) ## shuffle the deck then check the hand.
    shuffle!(deck)
    f(Hand(deck[1:5]))
  end
  @distributed (+) for i = 1:trials
    Int(checkHand(f))
  end  
end
\end{jlblock}
%
and running this:
%
\begin{jlblock}[parallel]
@time fh = paraCountHands(10_000_000,isFullHouse)
\end{jlblock}
returns \printpythontex[verbatim], which cuts the time by about quarter.

\subsection{parallelizing the ....}



\section{A Parallel Map Function}

\begin{jlcode}[parallel]
@everywhere function countHeads(n::Int)
   c::Int = 0
   for i=1:n
       c += rand(Bool)
   end
   c
end
\end{jlcode}

In Chapter \ref{ch:functional-programming}, we saw the \texttt{map} function, which is the standard way if you have an array and need to do the same thing to each element of the array, returning the result. If the function that is applied to each element of the array is costly and you have multiple processors/cores to work with, doing this in parallel can be helpful.

We demonstrate how to do this with just the coin flip right now using the \texttt{pmap} function.  We first make an array of length 12 (although this isn't an important number): 
%
\begin{jlblock}[parallel]
num_coins = 1_000_000_000*ones(Int64,12)
\end{jlblock}
%
so each element is 1 billion. We wish to run the coin flipper on each element of the array. Don't forget that the \texttt{countHeads} function must be declared everywhere. The following runs the function:
%
\begin{jlblock}[parallel]
@time pmap(countHeads,num_coins)
\end{jlblock}
%
and the resulting time is: \printpythontex[verbatim]

In comparison, if we run the regular map function:
%
\begin{jlblock}[parallel]
@time map(countHeads,num_coins)
\end{jlblock}
%
the result is \printpythontex[verb]. Again, this is about 4 times slower.


\subsection{When to use \texttt{pmap}?}

It seems like the \verb!pmap! function should always be used if it speeds up calculations.  However, note that in this example, the size of the array was small but the function on each element took a relatively long time.  If we have an array with millions of elements, though, the \texttt{pmap} function may actually be slower.  

\section{Shared arrays}

The last example shows that we may have to do something on an array. For
simplicity, let's say we have a fairly large array:

\begin{jlverbatim}
arr = rand(1:100,100_000_000);
\end{jlverbatim}

A common thing to do is to smooth an array (and is often done to images)
by a windowed mean, which means that every element is replaced by a mean
of the points around it in some window. We first define a windowed mean
by the following function:

\begin{jlverbatim}
function window_mean(arr::Array{Int64,1},i::Integer,width::Integer)
  window = max(1,i-width):min(i+width,length(arr))
  sum(arr[window])/(last(window)-first(window)+1)
end
\end{jlverbatim}

which first determine a window (being careful with the first and end of
the array) and then just calculating the mean.

Then if we have a new array of zeros:

\begin{jlverbatim}
smoothed_array = zeros(Float64,length(arr))
\end{jlverbatim}

we fill the new array with the windowed mean:

\begin{jlverbatim}
@time let
  for i=1:length(arr)
    smoothed_array[i]=window_mean(arr,i,100)
  end
end
\end{jlverbatim}

will fill the array with the smoothed version in about 43 seconds.

Note: the astute reader is probably thinking that using \texttt{map} to
do this is the right way to go, however, we can't use map in this
instance because the \texttt{window\_mean} function doesn't just use a
single number (from an array), it uses the entire array.

A natural way to speed this is is to send this to individual
processor/cores and work on this in pieces. One problem with this is
that we would have to send the entire array to each worker and that is
expensive. Since a 64-bit integer is 8 bytes, the array of 100 million
is about 800 Mb of memory and that is reasonable expensive to pass
around. Instead, we are going to use a shared array in the package
\texttt{SharedArrays} and you will need to add this.

\begin{jlverbatim}
using SharedArrays
\end{jlverbatim}
